% Cuerpo editor-ready del nuevo capitulo 57.

\section{Defecto suma--producto de APP y coeficiente lineal del producto modular de Euler}
\label{sec:l9s102:euler-defecto}

\subsection{Identidad exacta \texorpdfstring{\(459-405=54=[q]t_3(q)\)}{459-405=54=[q]t3(q)}}

Sobre el soporte de \(81\) celdas, las dos hojas de APP producen
\begin{equation}
 S_+=405,
 \qquad
 S_\times=459,
 \qquad
 S_\times-S_+=54.
 \label{eq:l9s102:euler-405-459}
\end{equation}
La regularización triádica suministra el exponente \(-1/12\). Su
publicación dodecafásica determina el polo \(q^{-1}\), y el cociente de
Dedekind de nivel \(3\) se escribe
\begin{equation}
 t_3(q)
 =\left(\frac{\eta(\tau)}{\eta(3\tau)}\right)^{12}
 =q^{-1}\prod_{n\ge1}(1+q^n+q^{2n})^{-12},
 \qquad q=e^{2\pi i\tau}.
 \label{eq:l9s102:euler-t3}
\end{equation}

\begin{theorem}[Identidad suma--producto de Euler--HMT]
\label{thm:l9s102:euler-54}
El defecto entre las dos evaluaciones de APP coincide con el coeficiente
lineal de \(t_3\):
\begin{equation}
 \boxed{
 S_\times-S_+=54=[q]\,t_3(q).}
 \label{eq:l9s102:euler-54}
\end{equation}
\end{theorem}

\begin{proof}
Cada factor admite
\((1+q^n+q^{2n})^{-12}=1-12q^n+O(q^{2n})\). Para el coeficiente de
\(q\) posterior al polo \(q^{-1}\), contribuyen las particiones de grado
\(2\) del producto de nivel \(3\). La expansión finita certificada produce
\([q]t_3=54\). APP obtiene el mismo entero por la resta exacta
\(459-405\).
\end{proof}

\subsection{Conservación del defecto mediante residuo, cociente y acarreo nonádico}

La identidad \cref{eq:l9s102:euler-54} enlaza la diferencia aditiva de dos
hojas finitas con un producto modular infinito. El enlace conserva el
origen de \(54\): suma, producto y regularización permanecen como tres
coordenadas distinguibles del estado.

\subsubsection{Completación EMRH y conservación de la unidad}
\label{sec:l9s102:euler-emrh}

La fibra ternaria de profundidad \(6\) posee \(3^6=729\) elementos. Su
completación decimal distingue
\begin{equation}
 729+270=999,
 \qquad
 729+271=1000,
 \qquad
 271=243+27+1.
 \label{eq:l9s102:euler-729-1000}
\end{equation}
El término terminal \(+1\) es la unidad de cierre transmitida por el
acarreo. La igualdad binomial
\begin{equation}
 1000=(9+1)^3=9^3+3\cdot9^2+3\cdot9+1
 \label{eq:l9s102:euler-binomial-1000}
\end{equation}
expresa el transporte entre la carta nonádica y la publicación decimal sin
perder la genealogía de la completación.

La evaluación arquimediana se realiza sobre cilindros semicerrados
compatibles. Las dos expansiones de un racional terminal representan el
mismo valor y conservan historias de frontera distintas. Esta estructura
tipa la identidad \(0.999\ldots=1\) como igualdad de evaluación acompañada
por memoria de ruta.

\section{Álgebras cuadráticas del TRIT y ley exponencial elíptica, parabólica e hiperbólica}
\label{sec:l9s102:euler-algebras-trit}

\subsection{Familia operatoria \texorpdfstring{\(J_{+1}^2=-I\), \(J_0^2=0\) y \(J_{-1}^2=I\)}{J+1^2=-I, J0^2=0 y J-1^2=I}}

Para \(\tau\in\{+1,0,-1\}\), definimos
\begin{equation}
 \mathbb A_\tau
 =\mathbb R[J_\tau]/(J_\tau^2+\tau I).
 \label{eq:l9s102:euler-algebras}
\end{equation}
La serie exponencial converge en cada álgebra y se reduce a
\begin{align}
 e^{tJ_{+1}}&=\cos t\,I+\sin t\,J_{+1},
 \label{eq:l9s102:euler-eliptica}\\
 e^{tJ_0}&=I+tJ_0,
 \label{eq:l9s102:euler-parabolica}\\
 e^{tJ_{-1}}&=\cosh t\,I+\sinh t\,J_{-1}.
 \label{eq:l9s102:euler-hiperbolica}
\end{align}
TRIT asigna respectivamente cierre elíptico, variación de primer orden y
apertura hiperbólica. La identidad clásica ocupa la hoja \(\tau=+1\):
\begin{equation}
 \boxed{\operatorname{Exp}(\pi J_{+1})+I=0.}
 \label{eq:l9s102:euler-clasica}
\end{equation}
La hoja parabólica conserva el primer jet, y la hoja hiperbólica publica
una pareja de autovalores recíprocos.

\subsection{Identidades \texorpdfstring{\(\operatorname{Exp}(\pi J_{+1})+I=0\)}{Exp(pi J+1)+I=0} y \texorpdfstring{\(\operatorname{Exp}(2\log\varphi\,H_\varphi)=F_{\mathrm{av}}^2\)}{Exp(2 log phi H_phi)=F_av^2}}

Sea
\begin{equation}
 F_{\mathrm{av}}=\begin{pmatrix}0&1\\1&1\end{pmatrix},
 \qquad
 H_\varphi=\frac{2F_{\mathrm{av}}^2-3I}{2\varphi-1}.
 \label{eq:l9s102:euler-hphi}
\end{equation}
Entonces \(H_\varphi^2=I\) y
\begin{equation}
 \boxed{
 \operatorname{Exp}(2\log\varphi\,H_\varphi)=F_{\mathrm{av}}^2.}
 \label{eq:l9s102:euler-fav}
\end{equation}
En esta coordinación, \(\pi\) mide el cierre angular, \(\varphi\) el
autovalor expansivo de la memoria y \(e\) la operación que transforma un
generador aditivo en transporte multiplicativo.

\section{Principio de Ambivalencia: involución entre \texorpdfstring{\(\pi/\varphi^2\)}{pi/phi^2} y \texorpdfstring{\(\varphi/\pi^2\)}{phi/pi^2} y acumulación nonádica de orientación}
\label{sec:l9s102:euler-ambivalencia}

\subsection{Acción \texorpdfstring{\(\mathcal Q(x,y)=(x/y^2,y/x^2)\)}{Q(x,y)=(x/y^2,y/x^2)} sobre los lectores areal y volumétrico}

Definimos el lector cuadrático
\begin{equation}
 \mathcal Q(x,y)
 =\left(\frac{x}{y^2},\frac{y}{x^2}\right).
 \label{eq:l9s102:euler-q}
\end{equation}
En las coordenadas
\begin{equation}
 P=xy,
 \qquad O=\frac{x}{y},
 \label{eq:l9s102:euler-po}
\end{equation}
su acción es
\begin{equation}
 P\circ\mathcal Q=P^{-1},
 \qquad
 O\circ\mathcal Q=O^3.
 \label{eq:l9s102:euler-q-action}
\end{equation}
La segunda iteración satisface
\begin{equation}
 \boxed{
 P\circ\mathcal Q^2=P,
 \qquad
 O\circ\mathcal Q^2=O^9.}
 \label{eq:l9s102:euler-q2}
\end{equation}
El producto retorna y la orientación conserva una potencia nonádica de
memoria.

\subsection{Identidad \texorpdfstring{\(\varphi^3(\pi/\varphi^2)^2(\varphi/\pi^2)=1\)}{phi^3(pi/phi^2)^2(phi/pi^2)=1} y bidireccionalidad tipada}

Para \((x,y)=(\pi,\varphi)\), las dos orientaciones son
\begin{equation}
 \rho_A=\frac{\pi}{\varphi^2},
 \qquad
 \rho_E=\frac{\varphi}{\pi^2},
 \label{eq:l9s102:euler-rho-a-e}
\end{equation}
y obedecen
\begin{equation}
 \boxed{\varphi^3\rho_A^2\rho_E=1.}
 \label{eq:l9s102:euler-ambivalencia-identidad}
\end{equation}
La involución \((x,y)\mapsto(y,x)\) intercambia las orientaciones areal y
electrónica del mismo lector \(R(x,y)=x/y^2\).

\section{Transformación exponencial del generador parangular \texorpdfstring{\(A_{\mathrm{rad}}I+C^*_{\mathrm{rad}}H_\varphi\)}{A_rad I+C*_rad H_phi}}
\label{sec:l9s102:euler-accion-par-angular}

La carta de acción
\begin{equation}
 T_{-34}(u)=u\,10^{-34}\mathcal U_S
 \label{eq:l9s102:euler-t34}
\end{equation}
transporta la identidad angular
\(C^*_{\deg}=2(\eta_{\mathrm{ret}}+\alpha_{\mathrm{HMT}})\) a
\begin{equation}
 \boxed{
 T_{-34}(C^*_{\deg})
 =2\hbar_{\mathrm{ret}}
 +2T_{-34}(\alpha_{\angle,\deg}).}
 \label{eq:l9s102:euler-cstar-hbar}
\end{equation}
Cada sumando pertenece así a la misma recta dimensional
\(\mathcal L_S\).

\subsection{Proyectores espectrales \texorpdfstring{\(P_\pm=(I\pm H_\varphi)/2\)}{P+/-=(I+/-H_phi)/2} y autovalores conjugados \texorpdfstring{\(q_\pm\)}{q+/-}}

Sea
\begin{equation}
 B=A_{\mathrm{rad}}I+C^*_{\mathrm{rad}}H_\varphi,
 \qquad
 P_\pm=\frac{I\pm H_\varphi}{2}.
 \label{eq:l9s102:euler-b}
\end{equation}
La exponenciación transforma el generador aditivo en dos hojas
multiplicativas conjugadas:
\begin{equation}
 \boxed{
 e^{-B}=q_+P_++q_-P_-,
 \qquad
 q_\pm=e^{-A_{\mathrm{rad}}\mp C^*_{\mathrm{rad}}}.}
 \label{eq:l9s102:euler-exp-b}
\end{equation}

\subsection{Compatibilidad con la publicación parangular y los canales \texorpdfstring{\(90/120\)}{90/120}}

Los cardinales \(90\) y \(120\) proceden de la incidencia trítica del TPK.
La pareja \(q_\pm\) no los genera: evalúa espectralmente el operador
\(V_{90,120}\) sobre esos canales previamente construidos.

\section{Frontera elíptica de acción e interior hiperbólico de Hilbert--Beltrami--Klein}
\label{sec:l9s102:euler-elipse-klein}

\subsection{Coordenadas conjugadas de la elipse de acción y ley \texorpdfstring{\(\operatorname{Area}(\theta)=\hbar_{\mathrm{ret}}\theta\)}{Area(theta)=hbar_ret theta}}

La rapidez angular
\begin{equation}
 \eta=\operatorname{artanh}\!\left(\frac{C^*}{A}\right)
 \label{eq:l9s102:euler-rapidez}
\end{equation}
define los semiejes simplécticos
\begin{equation}
 a_S=\sqrt{2\hbar_{\mathrm{ret}}e^\eta},
 \qquad
 b_S=\sqrt{2\hbar_{\mathrm{ret}}e^{-\eta}}.
 \label{eq:l9s102:euler-semiejes}
\end{equation}
Con
\begin{equation}
 Q=a_S\cos\theta,
 \qquad
 P=b_S\sin\theta,
 \label{eq:l9s102:euler-qp}
\end{equation}
la acción encerrada por el arco satisface
\begin{equation}
 \boxed{\operatorname{Area}(\theta)=\hbar_{\mathrm{ret}}\theta.}
 \label{eq:l9s102:euler-area-accion}
\end{equation}

\subsection{Correspondencia \texorpdfstring{\(\rho=\tan(\theta/2)=\tanh(u/2)\)}{rho=tan(theta/2)=tanh(u/2)} entre borde periódico e interior hiperbólico}

La coordenada de Gudermann
\begin{equation}
 \rho=\tan\frac\theta2=\tanh\frac u2
 \label{eq:l9s102:euler-gudermann}
\end{equation}
coordina las cartas circular e hiperbólica.
La frontera porta la órbita simpléctica; el dominio convexo interior porta
la métrica hiperbólica de Hilbert--Beltrami--Klein. La coordenada de
Gudermann realiza su correspondencia radial y conserva la distinción de
métricas.

\section{Compatibilidad del cierre operatorio de Euler con las publicaciones arquimediana y excepcional}
\label{sec:l9s102:euler-doble-proyeccion}

\subsection{Doble proyección del mismo estado enriquecido y conservación de la incidencia}

Sea \(X^{\mathrm{enr}}\) el estado enriquecido construido por
APP--TRIT--TPK. Sus dos publicaciones principales son
\begin{equation}
 X^{\mathrm{enr}}
 \xrightarrow{\ \mathcal P_{\mathrm{arq}}\ }\mathbb R,
 \qquad
 X^{\mathrm{enr}}
 \xrightarrow{\ \mathcal P_{\mathrm{exc}}\ }
 \mathrm{Witt}\to\mathrm{Golay}\to\mathrm{Leech}
 \to V^\natural\to\mathbb M\to J.
 \label{eq:l9s102:euler-doble-proyeccion}
\end{equation}
La primera contrae cilindros compatibles hacia el valor arquimediano; la
segunda conserva la incidencia de frontera hasta el corredor excepcional.
El entero \(54\) aparece en APP y en el coeficiente modular de
\cref{eq:l9s102:euler-54}; el corredor \(4\to2\to6\to54\) transporta
después esa incidencia hacia el retículo de Leech, el álgebra de operadores
de vértice y Moonshine. La acción del Monstruo reside en
\(V^\natural\); la genealogía digital y de rutas permanece en el dominio
del TPK y llega al corredor mediante los mapas de incidencia declarados.

\subsection{Sistema macrovectorial \texorpdfstring{\(\mathfrak E_{\mathrm{HMT}}=0\)}{E_HMT=0}: confluencia de las identidades eulerianas en un único operador}
\label{sec:l9s102:euler-macrovector}

Las identidades anteriores se reúnen en la aplicación tipada
\begin{equation}
 \boxed{
 \mathfrak E_{\mathrm{HMT}}=
 \begin{pmatrix}
 (S_\times-S_+)-[q]t_3\\[1mm]
 \operatorname{Exp}(I)-eI\\
 \operatorname{Exp}(\pi J_{+1})+I\\
 \operatorname{Exp}(J_0)-I-J_0\\
 \operatorname{Exp}(2\log\varphi\,H_\varphi)-F_{\mathrm{av}}^2\\
 \varphi^3\rho_A^2\rho_E-1\\
 T_{-34}(C^*)-2\hbar_{\mathrm{ret}}
 -2T_{-34}(\alpha_{\angle,\deg})
 \end{pmatrix}=0.}
 \label{eq:l9s102:euler-macrovector}
\end{equation}
Cada componente conserva su dominio y su demostración propia. La aplicación
agregada exhibe la coordinación de suma, producto, cierre elíptico, primer
jet, expansión hiperbólica, ambivalencia y acción.

La composición interna APP--TRIT--TPK ya transporta suma, producto,
orientación, residuo, cociente, acarreo, ruta y memoria hasta las
publicaciones anteriores. No se postula por ello un entrelazador HMT
ausente. Una promoción posterior sólo puede formularse mediante el cuadrado
tipado específico exigido por una representación ulterior ---por ejemplo,
una equivariancia grupal completa hasta Witt---, sin rebajar la confluencia
operatoria aquí demostrada.
